\section{Introduction}

Graph analysis has become fundamental to modern data science, powering applications ranging from recommendation systems in social networks and fraud detection in financial transactions to protein interaction analysis in bioinformatics and static code analysis.
As graph sizes grow to billions of vertices and edges, the need for efficient, scalable implementations becomes critical.
However, graph analysis algorithms are difficult to parallelize due to their irregular memory access patterns, poor data locality, and load imbalance across computational units.

%The GraphBLAS standard~\cite{9460572} addresses these challenges by providing a foundation for constructing high-performance parallel graph analysis algorithms.
%By expressing graph algorithms in the language of linear algebra~--- using generic sparse vectors and matrices as core data structures, together with a rich set of highly parameterizable operations (matrix multiplication, element-wise operations, etc.)~--- GraphBLAS offers a high-level, mathematically rigorous framework for algorithm construction.
%Operations over sparse generic matrices and vectors parameterized by semiring-like structures provide a powerful basis for expressing a wide range of graph analysis algorithms~\cite{8778338}.
%
%Despite its successes, GraphBLAS and its implementations face certain challenges including the implementation of high-performance generic operations over sparse data structures.
%Imbalanced subtasks together with irregular data access patterns make load balancing~--- and thus efficient parallel implementation~--- hard.
%Moreover, graph analysis algorithms are not single operations but nontrivial chains of operations over matrices and vectors.
%This leads to the problem of intermediate data structures.
%Kernel fusion could solve these problems, but it is too difficult to implement automatic optimizations for sparse generic computations, so it is done manually in the most frequent cases.

The GraphBLAS standard~\cite{9460572} addresses these challenges by providing a foundation for high-performance parallel graph analysis, expressing algorithms in linear algebra using generic sparse vectors/matrices and parameterizable operations. However, implementing high-performance generic operations over sparse data remains challenging: irregular access and imbalanced subtasks complicate load balancing and parallelization, while nontrivial chains of operations create intermediate data structures. Kernel fusion could help, but automatic optimization for sparse generic computations is difficult, so it is performed manually in the most common cases.

At the same time, Interaction Nets~\cite{10.1145/96709.96718}~--- a graph-rewriting model introduced by Yves Lafont~--- offer a computational model that is inherently well-suited for irregular parallelism.
In Interaction Nets, computation proceeds through local rewriting rules applied to pairs of connected nodes (so-called active pairs).
Crucially, multiple non-interfering active pairs can be reduced simultaneously, making Interaction Nets a promising foundation for parallel irregular computations, including graph analysis.
Furthermore, Interaction Nets are relatively close to functional programming languages, making it easy to transfer design approaches and ideas from the functional world.

However, comparisons between existing interaction-net-based evaluators remain very limited and are often performed only on basic benchmarks such as Fibonacci, Ackermann, sorting, and list traversal.

We believe that Interaction Nets can provide a powerful foundation for implementing graph analysis algorithms.
The long-term goal of our work is to investigate how Interaction Nets can help solve challenges with irregular parallelism and, possibly, with kernel fusion.
In this work, we do not focus on the GraphBLAS API implementation, but we are inspired by it: we implemented a minimal set of generic data structures and operations sufficient to implement several classic graph analysis algorithms.

This work is in progress, and our current contributions are as follows.
\begin{enumerate}
    \item We implemented a minimal, functionally equivalent subset of the functions over tree-represented matrices and vectors in F\# (a mature functional-first language, chosen for ease of prototyping and testing), Vine, and Inpla. 
    On this foundation, we implemented the following graph analysis algorithms: BFS, SSSP, and triangle counting.
    Thus, we enable direct comparison of the respective execution models and runtime systems.
    \item We evaluated the implemented algorithms on a set of real-world and synthetic graphs and highlight directions for future research. 
\end{enumerate}